\section{USAJMO 2015}
        \begin{enumerate}
        	\item (\textit{USAJMO 2015}) Given a sequence of real numbers, a move consists of choosing two terms and replacing each with their arithmetic mean. Show that there exists a sequence of 2015 distinct real numbers such that after one initial move is applied to the sequence -- no matter what move -- there is always a way to continue with a finite sequence of moves so as to obtain in the end a constant sequence.


 

	\item (\textit{USAJMO 2015}) Solve in integers the equation

 \[x^2+xy+y^2 = \left(\frac{x+y}{3}+1\right)^3.\]


 

	\item (\textit{USAJMO 2015}) Quadrilateral $APBQ$ is inscribed in circle $\omega$ with $\angle P = \angle Q = 90^{\circ}$ and $AP = AQ < BP$. Let $X$ be a variable point on segment $\overline{PQ}$. Line $AX$ meets $\omega$ again at $S$ (other than $A$). Point $T$ lies on arc $AQB$ of $\omega$ such that $\overline{XT}$ is perpendicular to $\overline{AX}$. Let $M$ denote the midpoint of chord $\overline{ST}$. As $X$ varies on segment $\overline{PQ}$, show that $M$ moves along a circle.


 

	\item (\textit{USAJMO 2015}) Find all functions $f:\mathbb{Q}\rightarrow\mathbb{Q}$ such that
 \[f(x)+f(t)=f(y)+f(z)\]for all rational numbers $x<y<z<t$ that form an arithmetic progression. ($\mathbb{Q}$ is the set of all rational numbers.)


 

	\item (\textit{USAJMO 2015}) Let $ABCD$ be a cyclic quadrilateral. Prove that there exists a point $X$ on segment $\overline{BD}$ such that $\angle BAC=\angle XAD$ and $\angle BCA=\angle XCD$ if and only if there exists a point $Y$ on segment $\overline{AC}$ such that $\angle CBD=\angle YBA$ and $\angle CDB=\angle YDA$.


 

	\item (\textit{USAJMO 2015}) Steve is piling $m\geq 1$ indistinguishable stones on the squares of an $n\times n$ grid. Each square can have an arbitrarily high pile of stones. After he finished piling his stones in some manner, he can then perform stone moves, defined as follows. Consider any four grid squares, which are corners of a rectangle, i.e. in positions $(i, k), (i, l), (j, k), (j, l)$ for some $1\leq i, j, k, l\leq n$, such that $i<j$ and $k<l$. A stone move consists of either removing one stone from each of $(i, k)$ and $(j, l)$ and moving them to $(i, l)$ and $(j, k)$ respectively, or removing one stone from each of $(i, l)$ and $(j, k)$ and moving them to $(i, k)$ and $(j, l)$ respectively.


 Two ways of piling the stones are equivalent if they can be obtained from one another by a sequence of stone moves.


 How many different non-equivalent ways can Steve pile the stones on the grid?


 

\end{enumerate}